\documentclass{article}
\usepackage{amsfonts, amsmath, amsthm, amssymb, enumerate, fancyhdr, gensymb, lastpage, mathtools, parskip, graphicx}
\usepackage{xcolor, tikz-cd}
\usepackage[a4paper, total={6in, 8in}]{geometry}
\usepackage{setspace}

\usepackage[OT2,T1]{fontenc}
\DeclareSymbolFont{cyrletters}{OT2}{wncyr}{m}{n}
\DeclareMathSymbol{\Sha}{\mathalpha}{cyrletters}{"58}

% Theorem environments
\newtheorem{theorem}{Theorem}
\newtheorem{lemma}[theorem]{Lemma}
\newtheorem{proposition}[theorem]{Proposition}
\newtheorem{corollary}[theorem]{Corollary}

% Definition-style environments
\theoremstyle{definition}
\newtheorem{definition}[theorem]{Definition}
\newtheorem{example}[theorem]{Example}

% Remark-style environments
\theoremstyle{remark}
\newtheorem{remark}[theorem]{Remark}


%\onehalfspacing
\setstretch{1.3}
\newcommand{\wg}[1]{\textcolor{violet}{#1}}

%\renewcommand{\C}{\mathcal C}
\newcommand{\Z}{\mathbb Z}
\newcommand{\Ztwo}{\mathbb Z / 2\mathbb Z}
\newcommand{\Ztwosq}{(\mathbb Z / 2 \mathbb Z)^2}
\newcommand{\GalK}{G_{\bar K / K}}
\newcommand{\Q}{\mathbb Q}
\newcommand{\R}{\mathbb R}
\newcommand{\im}{\text{im }}
\renewcommand{\div}{\, \vert \,}
\newcommand{\F}{\mathbb F}
\newcommand{\EQW}{E'(\Q)/2E(\Q)}
\newcommand{\Etors}{E(\Q)_{\text{tors}}}
\renewcommand{\F}[1]{\mathbb F_{#1}}
\newcommand{\Qst}{\mathbb Q(S,2)}
\newcommand{\Sphi}{S^\phi}
\newcommand{\Sphihat}{S^{\hat \phi}}
\newcommand{\bark}{{\bar{k}}}
\newcommand{\Qsqs}{(\Q^\times)/(\Q^\times)^2}
\newcommand{\Zn}[1]{\mathbb Z / #1 \mathbb Z}
\newcommand{\C}{\mathbb C}
\newcommand{\epq}{e^{p,q}}
\newcommand{\Spec}[1]{\text{Spec}(#1)}
\newcommand{\OX}{\mathcal O_X}

\newenvironment{wgenv}{\color{violet}}{}

\newcommand{\smalltwobytwo}[4]{
\left( \begin{smallmatrix*}
   #1 & #2 \\ #3 & #4 
\end{smallmatrix*}\right) 
}
\newcommand{\bmat}[1]{
    \begin{bmatrix*}
        #1
    \end{bmatrix*}
}

\title{
    Hodge Deligne Numbers for Hypersurfaces in Toric Varieties
}
\author{Will Gilroy}
\date{July 2026}

\begin{document}
\maketitle

\section{Mike Meeting 03 Aug}
Mike took us through some of the basic ideas. It was a long meeting,
but there were a lot of ideas, and so this is a writeup for some of
the ideas from that meeting.

Any \text{text in purple} indicates either a question from myself or
an incomplete thought. 

\subsection{Definition of Hodge-Deligne Numbers and Some Basic Properties}
Let $X$ be a variety over $\C$ of dimension\footnote{
    \wg{A variety is defined to be separated, irreducible, finite type
    $k$-scheme. I am trying to remember why we have that a variety has
    a well-defined dimension.}
}
 $n$.
The \textit{Hodge-Deligne numbers} denoted $e^{p,q}(X)$ for $0 \leq
p,q \leq n$ are a collection of integers satisfying properties soon to
be described. Moreover the polynomial $e_X(x,y) := \sum_{0 \leq p,q \leq n}
\epq(X) x^p y^q$ is called the \textit{Hodge-Deligne} polynomial.
The Hodge-Deligne numbers and polynomial satisfy the following
properties
\begin{itemize}
    \item The Hodge-Deligne polynomial is symmetric in $x,y$, $e_X(x,y) =
    e_X(y,x)$. In other words, the Hodge-Deligne numbers are symmetric
    in $p,q$, $e^{p,q}(X) = e^{q,p}$.
    
    \item If $X$ is reducible then the Hodge-Deligne polynomial
    decomposes additively over the components: if $X = Y \amalg Z$
    \wg{(presumably these have to be closed subsets?)} then $e_X = e_Y
    + e_Z$. 
    \wg{(A variety is by definition irreducible, so, perhaps the most
    general setting is for a separated, finite type over $\C$-scheme?
    Or maybe this property is supposed to be for closed subvarieties?)}

    \item If $X$ is a product variety then its Hodge-Deligne
    polynomial decomposes multiplicatively: if $X = Y \times_{\C} Z$
    then $e_X = e_Y e_Z$. 

    \item If $X$ is projective and smooth (over $\C$) then recall
    that the hodge numbers are the dimensions of the
    $\C$-vector spaces given by cohomology over the sheaf of
    differentials 
    \[
        h^{p,q}(X) = \dim_\C H^q(X, \Omega^p_X).
    \]

    In this setting, 
    the Hodge-Deligne numbers are related to the Hodge numbers in the
    following way
    \[
        \epq(X) = (-1)^{p+q} \cdot h^{p,q}(X). 
    \]
\end{itemize}

\subsection{Examples}

\begin{example}[A point]
    Let $X = \{*\} = \Spec \C$. We have that $X$ is smooth and
    projective. The following Hodge calculation will give that $e_X(x,y) = 1$.
    
    Since $X$ is affine $H^q(X, \Omega^p_X)$ is
    non-vanishing only for $q = 0$. 
    Recall that $\Omega^p_X = 0$ for all $p > 0$ \wg{(I have forgotten
    why this is true, but it follows because $X$ is a point. I think
    it should be easy to build the module of differentials in this
    case once one
    remembers the definition and construction)} and so $h^{p,q} = 0$ for all $p,q > 0$.
    The only possible non-zero Hodge number we may have is for $p = q=
    0$. We have the following
    \begin{align*}
        h^{0,0}(X) &:= H^0(X, \Omega_X^0) \\
        &= \dim_\C H^0(X, \OX) && \text{By definition}\\
        &= \dim_\C \Gamma(X, \OX) && \text{By the cohomology axioms} \\
        &= \dim_\C \C  && \text{$X = \Spec \C$ } \\ 
        &= 1.
    \end{align*}
    Overall, we have a that $h^{0,0}(X) = 1$ and 
    $h^{p,q}(X) = 0$ for all other $p,q$. 
    The claimed result follows.
\end{example}

\end{document}